\originalchapter{整体 Gauss--Bonnet 与 Gauss 映射}
\originalsection{有限三角剖分与曲率总量}
本章只引用拓扑中的两个工具: 紧致光滑曲面区域有有限光滑三角剖分, 可细分到每个三角形包含在一个坐标片内; 对任一此类剖分, $\chi=V-E+F$ 为 Euler 示性数, 不依赖剖分. 边界是有限条光滑闭曲线或分段光滑曲线, 三角剖分与其角点相容. 这些是拓扑工具, 不以曲率公式证明其存在与不变性.
\begin{theorem}[Gauss--Bonnet]\label{thm:global-gb}
设 $M$ 为紧致定向光滑曲面区域, 边界分段 $C^2$ 正则, 角点内角属于 $(0,2\pi)$, 且边界无自交. 边界按区域在左定向, 角点外角为 $\beta_j$. 则
\begin{equation}\label{eq:global-gb}
 \int_MK\dd A+\int_{\partial M}k_g\dd s+\sum_j\beta_j=2\pi\chi(M).
\end{equation}
无边界时即 $\int_MK\dd A=2\pi\chi(M)$.
\end{theorem}
\begin{proof}
取上述有限细三角剖分, 对每个三角形应用定理\ref{thm:local-gb}. 内部公共边由两个三角形沿相反方向走过, 测地曲率变号, 因而边积分相消; 边界边留下原边界积分. 记内部顶点数为 $V_i$, 边界顶点数为 $V_b$, 三角形数为 $F$. 每个三角形有3个角, 各外角为 $\pi$ 减该三角形内角, 所以把各局部公式相加得
\[
 \int_MK\dd A+\int_{\partial M}k_g\dd s
 =\sum_{\text{所有三角形角}}\alpha-\pi F.
\]
内部顶点周围各内角和为 $2\pi$. 光滑边界顶点周围为 $\pi$, 原有角点处则为其区域内角 $\alpha_j$. 在两侧加上原边界外角 $\pi-\alpha_j$, 得右侧总计 $2\pi V_i+\pi V_b-\pi F$.

记内部边数 $E_i$, 边界边数 $E_b$. 每条闭边界的边数等于顶点数, 故 $E_b=V_b$, 且三角形边计数给出 $3F=2E_i+E_b$. 因而
\[
 2\pi\chi=2\pi(V_i+V_b-E_i-E_b+F)
 =2\pi V_i+\pi V_b-\pi F,
\]
正是前面所得右侧. 无边界时取 $V_b=E_b=0$ 即得相应版本.
\end{proof}
\begin{example}
计算球面、环面以及有 $b$ 条边界的亏格 $g$ 连通定向紧致曲面区域的曲率约束.
\end{example}
\begin{solution}
球面 $\chi=2$, 所以总曲率为 $4\pi$; 半径改变使 $K$ 乘 $R^{-2}$、面积乘 $R^2$, 积分不变. 环面 $\chi=0$, 总曲率为0, 与第6章直接积分一致.

拓扑给出一般 $\chi=2-2g-b$. 若边界光滑且全为测地线, $k_g=0$, 则 $\int_MK\dd A=2\pi(2-2g-b)$. 例如带三个测地边界的亏格0区域, 即裤形区域, 总曲率为 $-2\pi$. 这说明在此条件下不可能处处 $K\ge0$.
\end{solution}
\begin{corollary}
紧致连通无边界定向曲面若 $K>0$, 则其亏格为0. 若 $K\le0$ 且不恒为0, 则亏格至少为2. 若亏格为1且 $K\ge0$ 或 $K\le0$, 则 $K\equiv0$.
\end{corollary}
\begin{proof}
用 $\int K=4\pi(1-g)$ 及连续函数积分的符号. 最后一项总积分为0, 不变号连续函数若不恒零会在某个开集上严格同号, 积分不为0, 矛盾.
\end{proof}
这里的亏格分类来自拓扑, 而曲率负责限制分类中哪些情况能实现. 对嵌入 $\R^3$ 的紧致曲面, 下节还会排除处处平直的可能.
\originalsection{Gauss 映射与椭圆点}
\begin{definition}
定向曲面的 Gauss 映射为 $N:M\to S^2$, 送每点到其单位法向. 目标球面按外法向定向.
\end{definition}
\begin{proposition}\label{prop:gauss-jacobian}
Gauss 映射的有向面积 Jacobian 为 $K$, 即 $N^*\dd A_{S^2}=K\dd A_M$. 绝对面积缩放为 $|K|$.
\end{proposition}
\begin{proof}
$p$ 与 $N(p)$ 的切平面均为 $N(p)^\perp$, 而正向切基以与 $N(p)$ 的叉积相容为条件, 两边可用同一正向正交基. $\dd N=-S$ 在二维上行列式为 $\det(-S)=\det S=K$. 面积换元的绝对值则为 $|K|$.
\end{proof}
这个等式解释了 $K$ 的“法向球面积密度”含义, 但在全局非单射时要计覆盖重数. 总有向面积有正负抵消, 不能直接等同于 Gauss 像点集的普通面积.
\begin{proposition}\label{prop:elliptic-point}
非空紧致无边界嵌入曲面至少有一点 $K>0$. 取原点在曲面之外并令 $|p|=R$ 在曲面达到最大, 在此处适当选外径向法向, 两主曲率都不大于 $-1/R$.
\end{proposition}
\begin{proof}
在最大点 $p$, 函数 $|X|^2$ 一阶导数为0, 所以 $p$ 垂直切平面, 可取 $N(p)=p/R$. 对任意单位切向 $V$, 取曲面曲线 $\gamma$ 经过 $p$ 且速度为 $V$. 二阶导数非正, 得
\[
 0\ge\tfrac12(|\gamma|^2)''=1+\langle p,\gamma''\rangle
 =1+R\II(V,V).
\]
因此所有单位方向的法曲率不大于 $-1/R$, 特别两个主曲率都严格为负, $K\ge1/R^2>0$. 反转法向仍保持 $K>0$.
\end{proof}
\begin{corollary}
任何紧致无边界嵌入环面都有 $K>0$ 的点和 $K<0$ 的点.
\end{corollary}
\begin{proof}
正值点由上命题得到. 总曲率为0, 连续性使正曲率在一个开集上贡献正积分, 故别处必须有负值.
\end{proof}
\originalsection{正曲率与整体凸性}
这里使用覆盖理论中的事实: 紧致连通空间到连通流形的局部微分同胚是有限覆盖; 连通覆盖的目标若为单连通球面, 则覆盖为一层. 在本册二维情形可直接说明第一点: 原像紧致而离散, 所以有限; 由逆函数定理在各原像点取互不相交小邻域, 再借紧致性排除其外的新原像, 得均匀覆盖邻域. 第二点是明确引用的拓扑结论.
\begin{theorem}[二维 Hadamard 凸性]\label{thm:hadamard}
紧致连通无边界嵌入曲面若 $K$ 处处非零, 则 $K>0$, Gauss 映射为微分同胚, 且每个切平面都是支撑平面. 它围成一个严格凸体.
\end{theorem}
\begin{proof}
由连续性和连通性, $K$ 不变号; 命题\ref{prop:elliptic-point}强迫其处处为正. 由于 $K>0$, 每点两种法向中恰有一种令形算子负定; 这个选择在局部光滑且在重叠片一致, 给出全局光滑法向. 取此法向, $\dd N=-S$ 处处可逆. Gauss 映射像是开集, 又因紧致为闭集, 所以覆盖整个 $S^2$. 由上述覆盖与单连通性, 它是一层覆盖, 即微分同胚.

主曲率乘积处处为正, 两者同号且不会越过0. 适当反转整体法向使最大距离点的形算子负定, 连通性使它处处负定. 对任意方向 $a\in S^2$, 高度函数 $h_a(p)=\langle p,a\rangle$ 的临界点恰为 $N(p)=a$ 或 $N(p)=-a$, 各有唯一一点. 在 $N=a$ 点, 高度 Hessian 为 $\II$, 负定, 是局部最大; 在 $N=-a$ 点它为 $-\II$, 正定, 是局部最小. 紧致性保证全局最大最小存在, 所以这两个点也是唯一全局最大和最小. 因而所有切平面都是支撑平面, 各方向最大支撑点唯一.

令 $C$ 为曲面的凸包. 它紧致: 三维空间中的凸组合总能缩到至多四个点, 因为超过四个点时仿射线性相关, 可沿相关系数调整权重直到至少一个权重为0; 重复即可. 因而凸包是紧集 $M^4$ 与紧致权重单纯形之连续像. 曲面不共面, 所以凸包有内部. 每个 $\partial C$ 点都有支撑平面: 取外部点趋近该边界点, 外部点到紧凸集的最近点给出支撑法向, 单位法向取收敛子列得到极限支撑面. 最近点性质的证明是对任意 $y\in C$ 比较到线段 $x+t(y-x)$ 的距离平方在 $t=0$ 的右导数, 得 $\langle a,y-x\rangle\le0$.

最大支撑值由曲面取得. 若凸组合落在对应支撑面上, 其所有正权重分量都必须达到该高度最大值, 而最大点唯一, 所以支撑面与 $C$ 的交仅有该点. 每个 $\partial C$ 点因此属于曲面. 反之每个曲面点都有支撑平面, 属于 $\partial C$, 故 $M=\partial C$. 其边界不含线段: 若含线段, 线段内点的支撑平面将同时包含两端点, 与单点支撑面矛盾. 所以 $C$ 是严格凸体.
\end{proof}
该定理关于紧致嵌入曲面. 不能把它用于开曲面或任意非单射浸入, 也不能由局部 $K>0$ 推断整曲面凸.
\originalsection{习题}
\begin{exercise}\label{ex:11a}
设紧致连通无边界定向曲面处处 $K=-1$, 亏格为 $g$. 求面积, 判断哪些 $g$ 从 Gauss--Bonnet 看可能. 说明该结论是否自动给出三维欧氏空间中的完整嵌入.
\end{exercise}
\begin{exercise}\label{ex:11b}
半径 $R$ 的球面上去掉两个互不相交极帽, 剩余纬带极角为 $\theta\in[\theta_1,\theta_2]$, $0<\theta_1<\theta_2<\pi$. 计算曲率积分和两条正向边界的测地曲率积分, 核对 Euler 示性数.
\end{exercise}
\begin{exercise}\label{ex:11c}
证明平面中光滑环形区域的两条正向边界测地曲率积分之和为0, 并解释为什么分别为 $2\pi$ 和 $-2\pi$.
\end{exercise}
